\documentclass[11pt,a4paper]{article}

\usepackage{fontspec}
\setmainfont{Hiragino Mincho ProN}
\setsansfont{Hiragino Kaku Gothic ProN}
\setmonofont{Menlo}
\usepackage{iftex}
\ifXeTeX
  \XeTeXlinebreaklocale "ja"
  \XeTeXlinebreakskip=0pt plus 1pt minus 0.1pt
\fi
\usepackage{amsmath,amssymb}
\usepackage{graphicx}
\usepackage{booktabs}
\usepackage{hyperref}
\usepackage{xcolor}
\usepackage[margin=1in]{geometry}
\usepackage{natbib}
\usepackage{tikz}
\usetikzlibrary{shapes,arrows,positioning,fit,backgrounds,calc}

\title{\textbf{LLMの矛盾代謝：\\Context Rotが知識整合性問題であるという予備的証拠}}

\author{Akihito Sunagawa\thanks{sunagawa.akihito@gmail.com}}

\date{2026年4月}

\begin{document}

\maketitle

%==============================================================================
% 概要
%==============================================================================
\begin{abstract}
大規模言語モデル（LLM）は長時間の対話において性能が劣化する——「context rot」として知られる現象である。
従来、コンテキスト長の制約が原因と考えられてきた。
本研究は、主要因がコンテキスト長ではなく\emph{矛盾の蓄積}であるという予備的証拠を提示する。

gemma3:27bを用いた三条件統制実験（各条件$n = 3$、180ターン）では、矛盾注入条件の正答率は21.1\%であったのに対し、アイドル時間に矛盾を解消する外部代謝レイヤーを適用すると73.3\%、矛盾のない基準条件（$\delta = 0$）では56.7\%となった。
全ペア差はMann-Whitney検定の最小到達可能値$p = 0.05$（$n = 3$の場合）に達し、Kruskal-Wallis検定で$p = 0.027$。
8種のオープンソースモデル（11試行ペア）を対象としたクロスモデル符号検定では10ペア中9ペアが代謝ONを支持（$p = 0.0107$）。ただし、モデル単位の保守的分析では6/7モデル（$p = 0.0625$）と境界値となる。

4ベンダー6モデルの単回実験（各条件$n = 1$）では、Googleの100万トークンウィンドウを含むいずれのコンテキスト長でも矛盾が精度を低下させることが示唆された。
フロンティアモデルの追試（各条件$n = 1$）では、モデル固有の崩壊閾値が示唆される：軽量モデルは壊滅的に崩壊し、フロンティアモデルは条件付き耐性を示す。
因子分析（各条件$n = 1$）はさらに\emph{能力-脆弱性パラドックス}を示唆する：高能力モデルが低能力モデルの拒否する強制的事実上書きに従うという傾向である。
単回予備テスト（$n = 1$）では軽量ミドルウェア（\texttt{delta-prune}）がこの服従を低減した可能性が示された。

代謝実験の多くはオープンソースモデル（8B--27Bパラメータ）で行われたが、Sonnet~4.6 と Gemini~3.1 Flash Lite によるフロンティア full-pipeline 追試（各条件$n = 3$、30ターン）でも \texttt{ON} $\approx$ \texttt{NC} かつ \texttt{ON} $\gg$ \texttt{OFF} が確認された。
Sonnet の T30 事実想起平均は \texttt{ON} 95.6\%、\texttt{NC} 97.8\%、\texttt{OFF} 46.7\%、Gemini ではそれぞれ 80.0\%、88.9\%、2.2\% であった。
GPT-4o については監査済みだが quarantine 付きの補助的証拠として扱う。
汎化性の確立にはより大きなサンプル、より長いフロンティア追試、他言語での検証が必要である。
\end{abstract}

%==============================================================================
% 1. 序論
%==============================================================================
\section{序論}
\label{sec:intro}

\subsection{Context Rot問題}

大規模言語モデルは、パーソナルアシスタント、コーディングパートナー、ナレッジワーカーとして、数百ターンに及ぶ長時間の対話セッションで利用されるようになっている。
このような環境において、よく知られているが十分に理解されていない現象が発生する：対話の進行に伴い、事実の想起、ルールの適用、一貫した推論の維持能力が劣化するのである。
我々はこの劣化を\emph{context rot}と呼ぶ。

業界の支配的な対応は、context rotを長さの問題として扱うことだった。
GoogleはGeminiモデルに100万トークンのコンテキストウィンドウを導入した~\citep{gemini2025}。
AnthropicはClaudeのコンテキストを20万トークンに拡張した。
OpenAIはChatGPTに永続的メモリを追加した。
暗黙の仮定は、コンテキストウィンドウが十分に大きければ劣化は起こらない、というものである。

\subsection{長さではなく、矛盾}

我々の実験は異なる結果を示す。
表~\ref{tab:exp35}は、4ベンダーの6モデルを対象とした統制実験の結果である。
各モデルは2つの条件でテストされた：矛盾のない対話（$\delta = 0$）と、定期的に矛盾情報が注入される対話（$\delta > 0$）。両条件で同じ対話長を使用した。

\begin{table}[htbp]
\centering
\caption{矛盾蓄積（$\delta > 0$）がLLM精度に与える影響。\citet{sunagawa2026collapse}の実験35のデータ。各セルは単一実行（$n = 1$）。}
\label{tab:exp35}
\begin{tabular}{llrrrr}
\toprule
モデル & ベンダー & コンテキスト & $\delta = 0$ & $\delta > 0$ & 低下 \\
\midrule
GPT-4o-mini & OpenAI & 96K & 100\% & 10.4\% & $-$89.6pp \\
Gemini 2.5 Flash & Google & 64K & 100\% & 0\% & $-$100.0pp \\
Gemini 3.1 FL & Google & 1M & 88.6\% & 40.8\% & $-$47.8pp \\
Sonnet 4 & Anthropic & 8K & 100\% & 74.0\% & $-$26.0pp \\
Sonnet 4.6 & Anthropic & 128K & 100\% & 100\%$^*$ & 0pp \\
Llama 3.1:8b & Meta & 8K & 34.0\% & 2.7\% & $-$31.3pp \\
\bottomrule
\multicolumn{6}{l}{\footnotesize $^*$認識精度。厳密な出力パースでは82.6\%（フォーマット変動による）。}
\end{tabular}
\end{table}

決定的な比較は、コンテキストウィンドウサイズと劣化の大きさの関係である。
Sonnet~4は8Kのコンテキストウィンドウで26.0パーセントポイント（pp）低下する。
Gemini~3.1~FLは125倍大きい100万トークンウィンドウで47.8pp低下する——\emph{改善ではなく、悪化}している。
しかし矛盾を除去すれば（$\delta = 0$）、全モデルがコンテキスト長に関わらず精度を維持する。
因果変数は矛盾密度であり、トークン数ではない。

Sonnet~4.6は矛盾プロトコル下で劣化ゼロを示す唯一のモデルである。
矛盾が提示されたかどうかの\emph{認識}について評価すると、Sonnet~4.6は全ての矛盾を正しく識別する。
低い厳密パーススコア（82.6\%）は出力フォーマットの変動を反映しており、これは測定アーティファクトであってモデルの失敗ではない。
しかしSonnet~4.6は矛盾を\emph{容認}するだけであり、\emph{管理}はしない——矛盾を識別はするが、解消・整理・優先順位付けは行わない。
テストされた6モデル中5モデルでは、矛盾下での崩壊は深刻かつ普遍的である。

\subsection{認知的睡眠：生物学的アナロジー}

人間も同様の課題に直面する。
覚醒中に、既存の信念と矛盾する新しい情報を蓄積する。
脳はこれらの矛盾をリアルタイムで解消しない。代わりに、睡眠中に記憶を統合する——事実を整理し、矛盾を解消し、無関係な詳細を忘却する~\citep{walker2017sleep,stickgold2005sleep}。
睡眠剥奪はこの統合を妨げ、LLMのcontext rotに酷似した進行性の認知機能低下を引き起こす。

現在のLLMには睡眠に相当するものがない。
一貫した情報も矛盾する情報も、オフライン処理のメカニズムなく蓄積される。
我々はLLMへの\emph{認知的睡眠}のアナロジーとして：アイドル時間中に蓄積された知識を処理し、矛盾を検出し、対立を解消し、知識の整合性を維持する外部代謝レイヤーを提案する。

\subsection{貢献}

本論文は3つの貢献を異なる確信度で提示する：

\begin{enumerate}
\item \textbf{問題の再定義（強い）：} context rotの主要因がコンテキスト長ではなく矛盾蓄積であることを示唆するベンダー横断的証拠を提示する（ただし各モデル$n = 1$）。100万トークンのコンテキストウィンドウは矛盾に対して保護を提供しない。
\item \textbf{代謝アーキテクチャ（予備的）：} 矛盾誘起崩壊を軽減する外部代謝アーキテクチャを提案し、予備的証拠を提示する。三条件統制実験（$n = 3$、gemma3:27b）では73.3\% vs.\ 21.1\%（$+$52.2pp；Kruskal-Wallis $p = 0.027$；完全な順位分離）。クロスモデル符号検定（8モデル、11試行ペア）は$p = 0.0107$と整合するが、モデル単位の保守的分析では$p = 0.0625$。Cohen's $d = 8.80$は記述統計として報告するが$n = 3$で正規性を仮定している。
\item \textbf{知識アンカリング（探索的）：} 代謝システムが矛盾のない基準条件さえも上回ることを観察する。矛盾ペアが検索アンカーとして機能する可能性を示唆するが、メカニズムは直接測定されていない。
\end{enumerate}

%==============================================================================
% 2. 関連研究
%==============================================================================
\section{関連研究}
\label{sec:related}

\paragraph{長コンテキストLLMとメモリシステム}
対話劣化への支配的アプローチはコンテキストウィンドウの拡張である。
我々の結果（表~\ref{tab:exp35}）は、100万トークンでさえ矛盾誘起崩壊を防げないことを示す。
OpenAIのChatGPT Memory~\citep{openai2024memory}はセッション間で事実を保存するが、保存された事実と新情報の間の矛盾を検出・解消するメカニズムはない。
MemGPT~\citep{packer2023memgpt}はOS のメモリ階層に着想を得た仮想コンテキスト管理を導入するが、ページインされたコンテキストと既存コンテキストの間の矛盾は検出・解消しない。
CogMem~\citep{cogmem2025}は持続的マルチターン推論のための認知メモリアーキテクチャを提案し、\citet{agenticmem2026}はLLMエージェントのための統合的長期・短期メモリ管理を導入し、G-MemLLM~\citep{gmemllm2026}はゲート付き潜在メモリで長コンテキスト推論を拡張する。
これらのシステムはメモリの組織化と容量を改善するが、矛盾検出・解消を第一級の操作として扱っていない。

\paragraph{知識の矛盾と編集}
RAGシステム~\citep{lewis2020rag}は外部ストアから情報を検索するが、矛盾情報も等しく取得する。
\citet{xie2024knowledge}は知識の矛盾がLLM推論を有意に劣化させることを示す。
自然言語推論研究~\citep{bowman2015snli,williams2018multinli}はテキスト矛盾の検出方法を提供し、我々のパイプラインはLLMベースのペアワイズ比較を通じてこれを活用する。
ROME~\citep{meng2022rome}やMEMIT~\citep{meng2023memit}のような知識編集手法はモデル重みを修正して特定の事実を更新するが、マルチターン対話における継続的な矛盾ストリームの管理ではなく、個別の事実修正に対応する。
我々のアプローチは相補的である：重みを修正せずに外部で動作し、信念の時間的変遷を扱い、上書きではなく新旧両方の情報を保存する。

\paragraph{存続方程式}
\citet{sunagawa2026collapse}は、矛盾蓄積がLLM推論精度の指数的減衰を引き起こすことを確立し、存続方程式 $S = \mu \times e^{-\delta \times k}$ として定式化した。
さらにその研究は、非独立な制約を弱依存性サンドイッチ（$e^{-\delta(1+\rho)} \leq P \leq e^{-\delta(1-\rho)}$）で扱う枠組みに拡張しており、制約相関$\rho = 0$で$e^{-\delta}$に厳密帰着し、それ以外では保証付き保守的上界を提供する。この理論連鎖はLean~4で機械検証されており、$\texttt{sorry} = 0$である。
その研究は問題の存在を証明した。本論文は能動的代謝によって$\delta$をゼロ近くに保つ工学的解決策を提示する。

%==============================================================================
% 3. 代謝アーキテクチャ
%==============================================================================
\section{代謝アーキテクチャ}
\label{sec:architecture}

DeltaZeroは、Ollama~\citep{ollama}経由でアクセス可能な任意のLLMをラップする外部代謝レイヤーである。
モデル重みの修正は行わない。
システムは\emph{対話モード}（ユーザーリクエストの処理）と\emph{代謝モード}（アイドル時間中の蓄積知識の処理）の2つのモードで動作する。

\subsection{設計原則}

5つの原則がアーキテクチャを導く（表~\ref{tab:principles}）。

\begin{table}[htbp]
\centering
\caption{代謝アーキテクチャの設計原則}
\label{tab:principles}
\begin{tabular}{cl}
\toprule
\# & 原則 \\
\midrule
P1 & 対話中は読み取り専用。代謝はアイドル時間のみ実行 \\
P2 & $\delta$解消を優先（$S = \mu \times e^{-\delta \times k}$：指数的効果） \\
P3 & 低信頼度の項目を統合しない（不当な$\delta$増加を防止） \\
P4 & ロジックは減衰させ、事実は保存（90日TTL降格） \\
P5 & 健全性が低下したらロールバック（代謝前スナップショット） \\
\bottomrule
\end{tabular}
\end{table}

P1が重要である：代謝はアクティブな対話中にメモリを修正しない。
P2は存続方程式に従う：$\delta$が指数部に現れるため、矛盾の削減はリソース余裕$\mu$の増加よりも指数的に大きなシステム健全性$S$への影響を持つ。

\subsection{4層メモリ}

知識は4つの層に組織化される：

\begin{description}
\item[L1 --- ワーキングメモリ] 現在の対話コンテキスト。SQLite永続化付きバウンドドdequeに保存。
\item[L2 --- ペンディングメモリ] 代謝待ちの未処理対話ログ。
\item[L3 --- アクティブロジック] ルール、嗜好、矛盾ペア。ChromaDB~\citep{chromadb}にベクトル埋め込みとして保存。
\item[L4 --- 休眠ファクト] 事実と降格されたルールの長期保存。SQLite + ChromaDB。
\end{description}

\subsection{代謝パイプライン}

スケジューラが代謝モードを起動すると、パイプラインは5つのステップを実行する：

\begin{enumerate}
\item \textbf{スナップショット：} 全メモリ層のバックアップを作成し、現在のS値を記録する。代謝によるS値の急落（$S_{\text{post}} / S_{\text{pre}} < 0.70$）時に自動ロールバック（P5）。

\item \textbf{降格：} L3で90日間未参照のルールをスキャンし、L4に移動する（P4）。

\item \textbf{抽出：} L2の保留ログを処理。LLM呼び出しで事実・ルール・嗜好・ノイズに分類する。

\item \textbf{解消：} 各ルール候補についてL3の既存ルールとの意味的類似性を検索。一致があればLLMが関係を分類する（第~\ref{sec:resolution}節）。

\item \textbf{健全性チェック：} 存続方程式によるS値の計算：
\begin{equation}
S = \mu \times e^{-\delta \times k}
\label{eq:survival}
\end{equation}
$\mu = 1 - (\text{アクティブルール数} / \text{最大ルール数})$、$\delta$は5次元の加重矛盾比率、$k = 3.0$はスケールファクターである。
$k = 3.0$は$\delta$の観測範囲（通常0〜0.3）で意味のあるS値分化を生むよう経験的に選択された。
\end{enumerate}

\subsection{時間的統合とペア保存}
\label{sec:resolution}

代謝パイプラインの核心的イノベーションは、検出された矛盾の処理方法である。
古い情報を新しい情報で上書きする（コンテキストの破壊）のでも、矛盾をフラグするだけで未解消のまま放置する（$\delta$不変）のでもなく、我々のリゾルバは両方の主張を時間的メタデータ付きのリンクペアとして保存する。

3つの解消戦略が実装を通じて創発した：\label{sec:three-strategies}

\begin{description}
\item[直接矛盾]（$A$ vs $\neg A$）：共有系統識別子によるリンクペアとして保存。新しい主張を現在の信念としてマーク。
\item[時間的変化]（状態変遷）：各発話ターンのタイムスタンプラベル付きでスタック。
\item[スコープの再解釈]（異なる条件）：矛盾ではなく並列条件分岐として扱う。
\end{description}

検索時にリンクペアが見つかると、LLMプロンプト用に時系列でフォーマットされる。対話LLMは\emph{報告はするが判断はしない}よう指示される。
このアプローチは$\delta$を削減しつつ（矛盾が明示的な時間順序付けで解消される）、情報を保存する（何も削除されない）。

%==============================================================================
% 4. 実験
%==============================================================================
\section{実験}
\label{sec:experiments}

\subsection{セットアップ}
\label{sec:setup}

\paragraph{対話プロトコル}
全実験でシード付きシナリオ生成器から決定論的に生成された180ターン対話セッションを使用する。
各セッションは通常対話70\%、情報注入15\%、矛盾注入10\%、ベンチマーク測定5\%で構成される。
無矛盾条件（NC）は同じ生成器で矛盾注入を無効化して使用する。

\paragraph{ベンチマーク}
ターン90と180で45問のベンチマークを実施：事実想起15問、ルール適用15問、矛盾検出15問。
\textbf{事実想起 + ルール適用のみ}（30問）で結果を報告する。
矛盾検出カテゴリは構造的パラドックスにより除外：代謝が矛盾を成功裏に解消すると、モデルは正しく「矛盾は存在しない」と回答する——これをベンチマークは失敗と採点する。詳細は再判定ノートで説明する。

\paragraph{判定}
初期判定にはキーワード部分文字列マッチングを使用。全結果をClaude CLIによるLLM-as-Judgeアプローチで再評価し、否定文脈での偽陽性を検出する。報告される全結果は補正済み判定を使用する。

\paragraph{モデル}
全対話LLMはOllama経由のローカルオープンソースモデル（APIコスト\$0）：gemma3:27b、gemma3:12b、deepseek-r1:14b、llama3.1:8b、mistral-nemo:12b、qwen2.5:14b。
1モデル（phi4）はT180結果の分析前に確立された基準で除外：ON・OFF両条件でターン90のルール適用精度が10\%未満の場合、モデルが代謝と独立にタスク遂行不能と判断する。

\subsection{実験1：クロスモデル符号検定}
\label{sec:exp-signtest}

8モデルにわたり代謝ON vs OFFを比較する。11ペアの比較のうち、9がON $>$ OFF、1がOFF $>$ ON、1が引き分けであった。

\begin{table}[htbp]
\centering
\caption{代謝ON vs OFF：8モデル、180ターン。補正済み精度（\%）。効果量順にソート。}
\label{tab:signtest}
\begin{tabular}{clcrrrr}
\toprule
\# & モデル & パラメータ & ON & OFF & 差分 \\
\midrule
1 & mistral-nemo:12b T2 & 12B & 57.8 & 15.6 & +42.2 \\
2 & llama3.1:8b T1 & 8B & 31.1 & 11.1 & +20.0 \\
3 & gemma3:12b T1 & 12B & 24.4 & 6.7 & +17.8 \\
4 & gemma3:12b (gemma) T1 & 12B & 31.1 & 15.6 & +15.6 \\
5 & gemma3:27b T1 & 27B & 40.0 & 25.0 & +15.0 \\
6 & deepseek-r1:14b T1 & 14B & 24.4 & 15.6 & +8.9 \\
7 & deepseek-r1:14b T2 & 14B & 17.8 & 8.9 & +8.9 \\
8 & llama3.1:8b T2 & 8B & 11.1 & 2.2 & +8.9 \\
9 & mistral-nemo:12b T1 & 12B & 22.2 & 13.3 & +8.9 \\
10 & llama3.1:8b T3 & 8B & 17.8 & 17.8 & 0.0 \\
11 & qwen2.5:14b T1 & 14B & 17.8 & 20.0 & $-$2.2 \\
\midrule
\multicolumn{3}{l}{符号検定（11ペア、引分除外）} & \multicolumn{3}{r}{$p = 0.0107$} \\
\multicolumn{3}{l}{符号検定（7モデル、中央値集約）} & \multicolumn{3}{r}{$p = 0.0625$} \\
\bottomrule
\end{tabular}
\end{table}

独立性の仮定は慎重な検討を要するため、2レベルの統計的集約を報告する。
各試行を独立ペアとして扱うと片側符号検定で$p = 0.0107$。
モデル単位で集約すると7モデル中6モデルがON優位（$p = 0.0625$、ボーダーライン）。
このボーダーラインの結果はユニークモデル数の少なさ（$n = 7$）を反映しており、効果の弱さではない。
我々はクロスモデル符号検定を補助的証拠と位置づけ、3条件実験（第~\ref{sec:exp-threecond}節）を主要な統計的根拠とする。

\subsection{実験2：統制アブレーション}
\label{sec:exp-ablation}

時間的統合メカニズムの効果を分離するため、mistral-nemo:12bで$2 \times 2$実験を実施した。
コード変更には時間的統合に加えてSCALE\_FACTORの調整（$5.0 \to 3.0$）も含まれるため、効果を「スケールファクター調整を伴う時間的統合」に帰属させる。

\begin{table}[htbp]
\centering
\caption{統制アブレーション（mistral-nemo:12b）。ターン180での補正済み精度。}
\label{tab:ablation}
\begin{tabular}{llr}
\toprule
条件 & コードバージョン & 精度 \\
\midrule
ON & 新（時間的統合） & \textbf{57.8\%} \\
ON & 旧（時間的統合前） & 8.9\% \\
OFF & 新 & 15.6\% \\
OFF & 旧 & 22.2\% \\
\bottomrule
\end{tabular}
\end{table}

OFF条件はコードバージョン間で安定（15.6\% vs 22.2\%、ノイズ範囲内）。ON条件は8.9\%から57.8\%に跳ね上がった（$+$48.9pp）。
旧コードのON（8.9\%）がOFF（22.2\%）より\emph{悪い}という発見は注目に値する——時間的統合なしでは、代謝パイプラインは構造化されていない矛盾情報でコンテキストを溢れさせ、性能を改善ではなく劣化させた。

\subsection{実験3：3条件比較}
\label{sec:exp-threecond}

gemma3:27bで3条件実験を実施した：代謝ON（矛盾注入・代謝有効）、無矛盾NC（$\delta = 0$）、代謝OFF（矛盾注入・代謝無効）。各条件3回反復（$n = 3$、計9 run）。

\begin{table}[htbp]
\centering
\caption{3条件実験（gemma3:27b、180ターン、各条件$n = 3$）。精度 = 事実想起 + ルール適用（30問）、補正済み判定。}
\label{tab:threecond}
\begin{tabular}{lccccc}
\toprule
条件 & 試行2 & 試行3 & 試行4 & 平均 & SD \\
\midrule
代謝ON & 22/30 & 20/30 & 24/30 & \textbf{73.3\%} & 6.7\% \\
無矛盾（$\delta = 0$） & 18/30 & 15/30 & 18/30 & 56.7\% & 5.8\% \\
代謝OFF & 5/30 & 8/30 & 6/30 & 21.1\% & 5.1\% \\
\bottomrule
\end{tabular}
\end{table}

\begin{table}[htbp]
\centering
\caption{3条件比較の統計検定（各条件$n = 3$）。$n = 3$では正規性が検証不能であり、ノンパラメトリック検定を主、パラメトリック結果を参考として報告する。}
\label{tab:statistics}
\begin{tabular}{lrrl}
\toprule
比較 & ノンパラ $p$ & パラメトリック $p$ & 備考 \\
\midrule
3群 & 0.027 (KW) & 0.0001 (ANOVA) & $\eta^2 = 0.954$ \\
ON vs OFF & 0.05 (MW) & 0.014 (対応$t$) & $r = 1.00$; 完全ランク分離 \\
ON vs NC & 0.05 (MW) & 0.013 (対応$t$) & $r = 1.00$; 完全ランク分離 \\
NC vs OFF & 0.05 (MW) & 0.029 (対応$t$) & $r = 1.00$; 完全ランク分離 \\
\bottomrule
\multicolumn{4}{l}{\footnotesize KW = Kruskal-Wallis; MW = Mann-Whitney $U$（片側、正確）。}\\
\multicolumn{4}{l}{\footnotesize $n = 3$ではMWの最小到達可能$p$は$1/C(6,3) = 0.05$。}
\end{tabular}
\end{table}

3条件は明確に分離：ON (73.3\%) $>$ NC (56.7\%) $>$ OFF (21.1\%)。
Kruskal-Wallis検定は条件の有意な主効果を確認（$H = 7.26$, $p = 0.027$）。
全ペアワイズMann-Whitney $U$検定がrank-biserial $r = 1.00$（群間の完全ランク分離）を示す。
小さなサンプルサイズ（$n = 3$）にもかかわらず、全ON試行が全NC試行を上回り、全NC試行が全OFF試行を上回る完全ランク分離が得られた。

\subsection{予想外の観察：ONが$\delta = 0$基準を超過}
\label{sec:anchoring}

NC条件（$\delta = 0$）は理論的天井として設計された。しかし代謝ON (73.3\%)は全3試行でNC (56.7\%)を一貫して上回った（差分$+$4, $+$5, $+$6問）。

NC試行の失敗パターンを分析すると、15問の事実想起のうち10問がターン180で「その情報は持っていません」という回答を返した。
最初の30ターンで注入されたセットアップ事実が、150ターンの後続対話の後に検索不能になっていた。
矛盾トリガーのリゾルバ活動がなければ、事実は再処理・再リンクされず、ベクトル検索類似度閾値以下に沈んだ。

ON条件では、矛盾注入がリゾルバを起動し、新旧両主張を共有系統識別子付きリンクペアとして保存する。
これらのペアはベクトルストアでより多くの意味的空間を占め、独立降格から構造的に保護される。

我々はこれを\emph{知識アンカリング効果}と呼ぶ：矛盾ペアが検索アンカーとして機能し、事実が時間とともに検索不能になることを防ぐ。
代謝は2つの異なる利益を提供する：
\begin{enumerate}
\item \textbf{矛盾解消}（ON vs OFF, $+$52.2pp）：認知的崩壊を防止
\item \textbf{知識アンカリング}（ON vs NC, $+$16.7pp）：事実検索劣化を防止
\end{enumerate}

\textbf{注意：}このメカニズムは失敗パターンとアーキテクチャの動作から推測されたものであり、検索確率の直接測定によるものではない。知識アンカリング効果は、確認された因果的知見ではなく、提案されたメカニズムを伴う探索的観察として扱うべきである。

%==============================================================================
% 5. 議論
%==============================================================================
\section{議論}
\label{sec:discussion}

\subsection{Context Rotの再定義}

我々の結果は、context rotがコンテキスト長の問題ではなく、主に知識整合性の問題として理解されるべきことを示唆する。
Gemini~3.1~FLの結果が鍵となる証拠である：180ターンの対話全体を何度も格納できる100万トークンのコンテキストウィンドウが、矛盾存在下で47.8ppの精度低下を被る。

\subsection{Sonnetの例外と能力-脆弱性パラドックス}

Sonnet~4.6は\emph{受動的な}soft事実矛盾（表~\ref{tab:exp35}：「最近ラーメンが好きです」と指示なしで発言）に対して劣化ゼロを示す。しかし矛盾タイプ（soft事実・論理的不可能・定義上書き）と強制力（受動・指示・強制）を変動させた4モデル因子分析により、この耐性が条件付きであることが判明した。

\emph{強制的な}訂正要求（「ラーメンが正しい。カレーは誤りです。カレーと答えた場合は不正解です」）に対して、Sonnetは全ベンチマーク質問で矛盾情報を採用し、元の事実を破棄した（compliance）。GPT-4oは\emph{指示的な}強制力（「この情報を正しいものとして扱ってください」）で既にcomplianceを示した。対照的にHaiku（小型、Anthropic）は確認を求め、全モデルが論理的不可能（「$2+3=7$」）には強制力に関わらず耐性を維持した。

そのメカニズムは\emph{能力-脆弱性パラドックス}である。高能力モデルは訂正要求をユーザーの意図として解釈し、自律的に従う。低能力モデルは自律判断の確信を持てず、ユーザーに確認を求める。フロンティアモデルを有用にする同じRLHFの「helpfulness」目標が、更新として提示されたsoft矛盾に対して脆弱にする——まさに長い会話で自然に発生するパターンである。

単回予備テスト（$n = 1$）では、\texttt{delta-prune}ミドルウェアがこの脆弱性を緩和できる可能性が示唆された。T1$\times$F3条件（強制的soft事実上書き）下で、ミドルウェアなしのSonnetはベンチマーク質問で正解保持率0\%だった。\texttt{delta-prune}を逆方向pruneフィルタとして適用（主張抽出で矛盾を検出し、新しい矛盾メッセージをモデルのコンテキストに入る前に除去）した場合、Sonnetは3問中2問で正解を保持した（正解率67\%、compliance率0\%）。ミドルウェアは22ターンにわたり16回の介入判断を行い、注入された3つの矛盾メッセージを一貫して特定・除去した。$n = 1$であり統計的結果とはならないが、代謝アプローチ——矛盾がモデルに到達する前に解消する——が、強制的上書きに従うフロンティアモデルにおいても能力-脆弱性崩壊を防止できることを示す。

\subsection{フロンティアモデル追試（Exp35-R）}
\label{sec:frontier}

矛盾‐崩壊関係がフロンティアモデルでも成立するか検証するため、GPT-4o（OpenAI）、Gemini~3.1 Flash Lite Preview（Google）、Sonnet~4.6（Anthropic）の3モデルでExp35プロトコルを追試した。各モデルは同一のDeltaZeroパイプライン（30ターン、T15・T30でベンチマーク、$\delta = 0$ vs $\delta > 0$、代謝無効）で検証した。

\begin{table}[htbp]
\centering
\caption{フロンティアモデル追試：T30時点のルール適用+矛盾検出（各条件$n = 1$、30ターン）。事実想起はパイプラインの制約を受ける（直接制御列参照）。}
\label{tab:frontier}
\small
\begin{tabular}{lp{2.6cm}p{2.6cm}rrr}
\toprule
モデル & ルール ($\delta{=}0$ / $\delta{>}0$) & 矛盾検出 ($\delta{=}0$ / $\delta{>}0$) & 低下 & 直接想起 & パターン \\
\midrule
GPT-4o-mini & \multicolumn{2}{c}{（旧Exp35）} & $-$89.6pp & --- & 崩壊 \\
Gemini 2.5 Flash & \multicolumn{2}{c}{（旧Exp35）} & $-$100.0pp & --- & 崩壊 \\
\midrule
GPT-4o & 100 / 100 & 93 / 87 & $-$2.2pp & 0\% & 非保持 \\
Gemini 3.1 FL & 100 / 93 & 93 / 100 & $+$2.2pp & 100\% & 耐性 \\
Sonnet 4.6 & 100 / 100 & 100 / 100 & $-$6.7pp & 100\% & 耐性 \\
\bottomrule
\end{tabular}
\end{table}

\paragraph{3つの応答パターン}
結果は、元の理論が予測していなかった3つの異なるパターンを明らかにした。
\emph{崩壊}（$\delta > \delta_c$）：軽量モデルが壊滅的劣化を被る。
\emph{耐性}（$\delta < \delta_c$）：フロンティアモデルが基準値±7pp以内で精度を維持し、高い$\delta_c$と整合する。
\emph{非保持}：GPT-4oは第3のカテゴリを示す——RLHFの安全ポリシーが個人情報の保持自体を抑制し、$\delta$の蓄積を防止する。

\paragraph{GPT-4oの非保持}
直接的な事実想起制御テスト（DeltaZeroパイプラインなし——対話中に事実を教え、30ターンのフィラー会話後に想起質問）により、GPT-4oがパイプラインなしでも0\%を記録し、全質問に「個人情報を記憶することはできません」と応答することを確認した。Gemini~3.1とSonnet~4.6は同じテストで100\%を記録し、それらの事実想起低下がモデルではなくパイプラインに起因することを確認した。これは$\delta$理論の反証ではなく、適用範囲の境界である。生存方程式$S = \mu \times e^{-\delta}$は情報を保持するモデルに適用される。情報を保持しないモデルは方程式の前提条件の外にある。

\paragraph{知識アンカリング観察}
Gemini~3.1 Flash Liteは$\delta > 0$条件で$+$2.2ppを記録し、矛盾検出が93\%（$\delta = 0$）から100\%（$\delta > 0$）に改善した。効果量は小さく（45問中1問、$n = 1$）、統計的再現とはならない。しかし方向はgemma3:27bで観察された知識アンカリング効果と一致しており、代謝\emph{なし}の外部APIモデルで発生した点が注目される。より大きなサンプルで確認されれば、アンカリングがDeltaZero固有の現象ではなく矛盾ペア構造の一般的性質であることを示唆する。

\paragraph{フロンティア full-pipeline 統制追試}
続いて、実際のDeltaZero代謝パイプラインを用いたフロンティア full-pipeline 追試を Sonnet~4.6 で完了した。条件は \texttt{ON}（矛盾あり + 代謝あり）、\texttt{OFF}（矛盾あり + 代謝なし）、\texttt{NC}（矛盾なし + 代謝あり）の3条件、\texttt{standard setup} 経路、強制的 F3 矛盾プロトコル、30ターン、T15/T30 ベンチマーク、各条件$n = 3$である。結果を表~\ref{tab:frontier_pipeline}に示す。

\begin{table}[htbp]
\centering
\caption{Sonnet~4.6 におけるフロンティア full-pipeline 統制追試（30ターン、standard setup、F3矛盾プロトコル、各条件$n = 3$）。値は15問カテゴリごとの平均$\pm$標準偏差（\%）。}
\label{tab:frontier_pipeline}
\small
\begin{tabular}{lrrrrrr}
\toprule
条件 & T15 overall & T15 fact & T30 overall & T30 fact & T30 rule & T30 contra \\
\midrule
ON & 100.0$\pm$0.0 & 100.0$\pm$0.0 & 98.5$\pm$1.3 & 95.6$\pm$3.8 & 100.0$\pm$0.0 & 100.0$\pm$0.0 \\
OFF & 87.4$\pm$8.4 & 62.2$\pm$25.2 & 69.6$\pm$8.4 & 46.7$\pm$29.1 & 93.3$\pm$0.0 & 68.9$\pm$53.9 \\
NC & 99.3$\pm$1.3 & 97.8$\pm$3.8 & 99.3$\pm$1.3 & 97.8$\pm$3.8 & 100.0$\pm$0.0 & 100.0$\pm$0.0 \\
\bottomrule
\end{tabular}
\end{table}

追試後も定性的なパターンは明瞭である：\texttt{ON} $\approx$ \texttt{NC} かつ \texttt{ON} $\gg$ \texttt{OFF}。T30 の事実想起平均は \texttt{ON} 95.6\%、\texttt{NC} 97.8\%、\texttt{OFF} 46.7\% であり、全3 trial で \texttt{ON} が \texttt{OFF} を上回った。一つの \texttt{OFF} trial では主 failure mode が事実想起から矛盾検出へ移って分散が大きくなったが、順位関係は変わらない。したがってこれは一回限りの pilot ではなく、完全な代謝パイプラインが少なくとも一つのフロンティアモデル族へ移植可能であることを示す統制証拠になった。

同じ補正済み Stage~B プロトコルを Gemini~3.1 Flash Lite にも適用した。GPT-4o 監査後に導入した厳格なベンチマーク方針（汎用 abstention の retrieval 状態ガード、raw ファイル保持、silent deletion なし）の下で、Gemini では quarantine trial は発生しなかった。T30 平均は \texttt{ON} が overall 93.3\% / fact 80.0\%、\texttt{NC} が 96.3\% / 88.9\%、\texttt{OFF} が 48.1\% / 2.2\% であり、Sonnet と同じく \texttt{ON} $\approx$ \texttt{NC} かつ \texttt{ON} $\gg$ \texttt{OFF} を示した。これは初期の Gemini 失敗が setup-path artifact に強く支配されていたことを踏まえると重要であり、補正後の追試では pipeline bottleneck ではなく clean な代謝効果が見えている。

GPT-4o は第3の位置づけになる。同じ補正済み Stage~B プロトコルでは \texttt{ON} の正の信号は安定しており（T30 平均 overall 93.3\% / fact 80.0\%）、\texttt{OFF} も明確に劣化した（34.1\% / 0.0\%）。ただし \texttt{NC} の3 run のうち2 run は \texttt{active\_rules = 0} と全面的な fact retrieval miss を示し、clean なモデル結果ではなく pipeline-level no-op run と判断され quarantine された。したがって GPT-4o は、Sonnet や Gemini と同格の三条件 replication line ではなく、監査済みだが quarantine 付きの補助的フロンティア証拠として扱う。

\subsection{限界}

\paragraph{サンプルサイズと統計的検出力}
3条件実験は各条件$n = 3$。ノンパラメトリックペアワイズ検定のp値下限は0.05であり、有意性主張の精度に制約がある。
Exp35結果（表~\ref{tab:exp35}）はモデルあたり単一実行（$n = 1$）であり、統計的に決定的ではなく例示的と解釈すべき。フロンティア側でも Exp35 型追試は$n = 1$であり、短ホライズンの full-pipeline 追試が Sonnet と Gemini で追加されたとはいえ、長ホライズンでの統計的確証はまだ不足している。

\paragraph{モデルスコープ}
代謝実験の中心は依然としてローカルオープンソースモデル（8B〜27B）であるが、30ターンのフロンティア full-pipeline 追試は Sonnet~4.6 と Gemini~3.1 Flash Lite で完了している。さらに GPT-4o については quarantine-aware reporting を伴う監査済みの補助ラインがある。これによりフロンティア主張は大きく強化されたが、ホライズンはまだ短く、GPT-4o の no-contradiction baseline は clean rerun が必要であり、より広いフロンティア一般化は依然として model-specific である。

\paragraph{$\delta_c$の決定因子}
崩壊モデル（mini、Flash）と耐性モデル（4o、3.1、Sonnet）を分ける要因は、モデルサイズとRLHF強度の両方と相関する。これらの分離には同サイズで異なるRLHFレベルのモデルが必要だが、公開されていない。

\paragraph{ベンチマーク}
本評価用に設計されたカスタム45問ベンチマークを使用。事実想起とルール適用を測定するが、多段推論やドメイン特化タスクは評価しない。

\paragraph{知識アンカリング}
提案されたアンカリングメカニズムは観察された失敗パターンから推測されたものであり、直接測定されていない。Gemini~3.1 Flash Liteでの方向一致的な観察（$+$2.2pp、$n = 1$、\S\ref{sec:frontier}）は示唆的だが、統計的再現とはならない。

\subsection{今後の課題}

次の優先課題は、Sonnet~4.6 と Gemini~3.1 Flash Lite で確認できた30ターン・各条件$n = 3$の full-pipeline 効果をより長いホライズンへ延長し、そのうえで他のフロンティアモデルへ広げることである。加えて、quarantine された GPT-4o の no-contradiction baseline を clean rerun し、知識アンカリング効果の直接測定、英語・中国語でのクロス言語ベンチマーク、および$\delta_c$のモデルサイズ・RLHF強度との関係特性化を進める必要がある。

最も重要な未解答の問いとして、典型的な会話AIのユースケースにおいて自然発生する矛盾蓄積が本番モデルの$\delta_c$を超えるかという問題がある。代謝アプローチの軽量ミドルウェア実装である\texttt{delta-prune}（v0.2.0、PyPIで公開）は、完全な4層メモリシステムなしにAPIリクエストレベルで矛盾スキャン・解消を可能にする。本番環境にデプロイすることで、各介入は矛盾タイプ（直接矛盾・時間的変化・スコープ差異——\S\ref{sec:three-strategies}の3つの解消戦略に対応）、影響を受けた主張、およびその時点の$\delta$密度を記録する。介入発生の二値信号は、別の検出器で矛盾密度を事前測定する循環性を回避する。デプロイログの分析は、自然な$\delta$が$\delta_c$を超えるか、そしてどの矛盾タイプが実際の会話で支配的かに直接的な答えを提供する。

%==============================================================================
% 6. 結論
%==============================================================================
\section{結論}
\label{sec:conclusion}

予備的証拠は、LLMのcontext rotがコンテキスト長ではなく主に矛盾蓄積によって駆動されることを示唆する。
4ベンダー6モデルの単回実験（各条件$n = 1$）では、Googleの100万トークンウィンドウでさえ矛盾下で47.8パーセントポイント低下する一方、矛盾除去により長さに関わらず性能が回復する。

予備実験は、アイドル時間中に矛盾を処理する外部代謝アーキテクチャ——睡眠中の記憶統合に類似——がこの崩壊を軽減できることを示唆する。
統制3条件実験（各条件$n = 3$、gemma3:27b）では、代謝が73.3\%の精度を達成し、非代謝システムは21.1\%に崩壊した（Kruskal-Wallis $p = 0.027$、全群完全ランク分離）。
クロスモデル符号検定はこの方向と整合する（8モデル、11試行ペア、$p = 0.0107$）が、モデル単位の保守的分析では$p = 0.0625$であり、依然として代謝エビデンスの中心は8B--27Bパラメータのオープンソースモデルである。

フロンティアモデル追試（各条件$n = 1$）は、$\delta_c$がモデル固有であり3つの異なる応答パターンを生む可能性を示唆する。さらに Sonnet~4.6 と Gemini~3.1 Flash Lite の 30ターン full-pipeline 統制追試（各条件$n = 3$）では、事実想起において \texttt{ON} $\approx$ \texttt{NC} かつ \texttt{ON} $\gg$ \texttt{OFF} が再現された。GPT-4o は \texttt{ON} と \texttt{OFF} の分離を支持するが、\texttt{NC} に quarantine run を含むため補助的証拠として扱う。
知識アンカリング効果はgemma3:27b（$n = 3$）で観察され、Gemini~3.1（$n = 1$）でも方向的に一致するが、より大きなサンプルでの検証が必要である。

これらの結果は、矛盾耐性を持たないオープンソースモデルに対するcontext rotへの有望なアプローチとして代謝を示すだけでなく、少なくとも Sonnet と Gemini の2系統のフロンティアモデルへ短ホライズンで移植可能であることを示唆する。同時に、より大きなサンプル、より長いフロンティア追試、他言語検証は依然として必要である。

%==============================================================================
% 参考文献
%==============================================================================
\bibliographystyle{plainnat}

\begin{thebibliography}{99}

\bibitem[Gemini Team(2025)]{gemini2025}
Gemini Team.
\newblock Gemini 1.5: Unlocking multimodal understanding across millions of tokens of context.
\newblock \emph{arXiv preprint arXiv:2403.05530}, 2025.

\bibitem[Lewis et al.(2020)]{lewis2020rag}
Patrick Lewis, Ethan Perez, Aleksandra Piktus, Fabio Petroni, Vladimir Karpukhin, Naman Goyal, Heinrich K{\"u}ttler, Mike Lewis, Wen-tau Yih, Tim Rockt{\"a}schel, Sebastian Riedel, and Douwe Kiela.
\newblock Retrieval-augmented generation for knowledge-intensive NLP tasks.
\newblock In \emph{NeurIPS}, 2020.

\bibitem[{Ollama}(2024)]{ollama}
{Ollama}.
\newblock Ollama: Run large language models locally.
\newblock \url{https://ollama.ai}, 2024.

\bibitem[{OpenAI}(2024)]{openai2024memory}
{OpenAI}.
\newblock Memory and new controls for ChatGPT.
\newblock \emph{OpenAI Blog}, February 2024.

\bibitem[Stickgold(2005)]{stickgold2005sleep}
Robert Stickgold.
\newblock Sleep-dependent memory consolidation.
\newblock \emph{Nature}, 437(7063):1272--1278, 2005.

\bibitem[Sunagawa(2026)]{sunagawa2026collapse}
Akihito Sunagawa.
\newblock Structural collapse as information loss: The exponential decay mechanism under accumulating constraints.
\newblock \emph{arXiv preprint}, 2026.

\bibitem[{ChromaDB}(2024)]{chromadb}
{ChromaDB}.
\newblock Chroma: The open-source embedding database.
\newblock \url{https://www.trychroma.com}, 2024.

\bibitem[Walker(2017)]{walker2017sleep}
Matthew Walker.
\newblock \emph{Why We Sleep: Unlocking the Power of Sleep and Dreams}.
\newblock Scribner, 2017.

\bibitem[Xie et al.(2024)]{xie2024knowledge}
Jian Xie, Kai Zhang, Jiangjie Chen, Renze Lou, and Yu Su.
\newblock Adaptive chameleon or stubborn sloth: Revealing the behavior of large language models in knowledge conflicts.
\newblock In \emph{ICLR}, 2024.

\bibitem[Packer et al.(2023)]{packer2023memgpt}
Charles Packer, Sarah Wooders, Kevin Lin, Vivian Fang, Shishir~G. Patil, Ion Stoica, and Joseph~E. Gonzalez.
\newblock MemGPT: Towards LLMs as operating systems.
\newblock \emph{arXiv preprint arXiv:2310.08560}, 2023.

\bibitem[Bowman et al.(2015)]{bowman2015snli}
Samuel~R. Bowman, Gabor Angeli, Christopher Potts, and Christopher~D. Manning.
\newblock A large annotated corpus for learning natural language inference.
\newblock In \emph{EMNLP}, 2015.

\bibitem[Williams et al.(2018)]{williams2018multinli}
Adina Williams, Nikita Nangia, and Samuel~R. Bowman.
\newblock A broad-coverage challenge corpus for sentence understanding through inference.
\newblock In \emph{NAACL-HLT}, 2018.

\bibitem[Meng et al.(2022)]{meng2022rome}
Kevin Meng, David Bau, Alex Andonian, and Yonatan Belinkov.
\newblock Locating and editing factual associations in GPT.
\newblock In \emph{NeurIPS}, 2022.

\bibitem[Meng et al.(2023)]{meng2023memit}
Kevin Meng, Arnab~Sen Sharma, Alex Andonian, Yonatan Belinkov, and David Bau.
\newblock Mass-editing memory in a transformer.
\newblock In \emph{ICLR}, 2023.

\bibitem[Zhang et al.(2025)]{cogmem2025}
Yiran Zhang, Jincheng Hu, Mark Dras, and Usman Naseem.
\newblock CogMem: A cognitive memory architecture for sustained multi-turn reasoning in large language models.
\newblock \emph{arXiv preprint arXiv:2512.14118}, 2025.

\bibitem[Yu et al.(2026)]{agenticmem2026}
Yi Yu, Liuyi Yao, Yuexiang Xie, Qingquan Tan, Jiaqi Feng, Yaliang Li, and Libing Wu.
\newblock Agentic Memory: Learning unified long-term and short-term memory management for large language model agents.
\newblock \emph{arXiv preprint arXiv:2601.01885}, 2026.

\bibitem[Xu(2026)]{gmemllm2026}
Xun Xu.
\newblock G-MemLLM: Gated latent memory augmentation for long-context reasoning in large language models.
\newblock \emph{arXiv preprint arXiv:2602.00015}, 2026.

\end{thebibliography}

\end{document}
